\documentclass[11pt]{article}
\usepackage[utf8]{inputenc}
\usepackage{amsmath,amsfonts,amssymb,amsthm}
\usepackage{mathtools}
\usepackage{geometry}
\usepackage{fancyhdr}
\usepackage{xcolor}
\usepackage{enumitem}
\usepackage{tcolorbox}
\usepackage{physics}
\usepackage{microtype}
\geometry{margin=2.5cm}

\title{\Huge Foundations of Optimal Transport}
\author{Pierre Chambet}
\date{May 2025}

\pagestyle{fancy}
\fancyhf{}
\rhead{Optimal Transport Foundations}
\lhead{Level 0}
\cfoot{\thepage}

% Highlight box
\newtcolorbox{highlight}{colback=blue!5!white,colframe=blue!75!black,boxrule=0.8pt,arc=4pt,left=6pt,right=6pt,top=6pt,bottom=6pt}

\begin{document}

\maketitle

\section*{About This Document – \textit{Optimal Transport Foundations}}

\vspace{1em}

This document is the mathematical companion to the slides I presented at the workshop of the \textbf{Computational Systems Biology Laboratory} at NAIST (Nara Institute of Science and Technology) in May 2025. It covers the main ideas of optimal transport rather than every detail, and assumes some familiarity with probability and analysis.

\vspace{1em}


You’ll find:
\begin{itemize}
  \item A visual and intuitive introduction to transporting mass;
  \item A clear recap on probability measures, Wasserstein distances, and duality;
  \item An explanation of both discrete and continuous cases (Monge vs Kantorovich);
  \item And finally, a bridge to generative models that leverage Optimal Transport.
\end{itemize}

\vspace{1em}

This foundational summary also serves as a stepping stone toward deeper theoretical exploration.  
In particular, I will build upon this work to analyze and understand the recent paper:

\begin{center}
\textbf{\textit{“Generative Modeling with Optimal Transport Maps”}} (ICLR 2022)\\
by Litu Rout, Evgeny Burnaev, and Alexander Korotin.
\end{center}

\vspace{1em}


\begin{highlight}
    \begin{center}
       \textbf{Let's start} 
    \end{center}
\end{highlight}



\newpage
\lhead{Level 0}
\section*{Level 0 - Introduction: The Big Picture}
Imagine you have a pile of sand arranged in one shape on a table---perhaps a tall, pointy pyramid. Now, you want to rearrange that same sand into another shape---say, a smooth, round mound. The total amount of sand remains the same.

\vspace{0.5em}
\noindent This gives rise to a natural and powerful question:

\begin{center}
\textit{What is the most efficient way to move the sand from the source shape to the target shape?}
\end{center}

This is the central idea behind \textbf{Optimal Transport (OT)}: it provides a mathematically rigorous way to compare two distributions by asking how much ``effort'' it takes to morph one into the other.

Before diving into more advanced concepts like Wasserstein distances or generative models, we need to build a rock-solid understanding of the foundational ideas. That's our goal here.

\section*{1. Probability Measures: The Language of Mass}

Let $X$ be a space (like $\mathbb{R}^n$), and let $\mu$ be a \textbf{probability measure} on $X$.

\begin{itemize}
    \item Think of $\mu$ as describing how mass or probability is spread over $X$.
    \item If $A \subset X$ is a subset, then $\mu(A)$ tells you how much mass is located in $A$.
    \item In simple terms: $\mu$ is a generalization of histograms to continuous spaces.
\end{itemize}

\textbf{Example (Discrete case):} In $\mathbb{R}^2$, suppose:
\[
\mu = \frac{1}{3}\delta_{x_1} + \frac{1}{3}\delta_{x_2} + \frac{1}{3}\delta_{x_3}
\]
This describes three equally weighted particles at positions $x_1$, $x_2$, $x_3$.

\section*{2. Transporting a Measure: Pushforward}

Now, suppose we apply a function $T : X \to Y$ that moves every point $x$ in $X$ to some point $T(x)$ in $Y$.

This induces a new distribution $T_\#\mu$ on $Y$, called the \textbf{pushforward of $\mu$ by $T$}. Formally:
\[
T_\#\mu(B) := \mu\left(T^{-1}(B)\right) \quad \text{for all measurable } B \subset Y
\]

\begin{highlight}
\textit{This means that the amount of mass that ends up in region $B$ of $Y$ is exactly the amount of mass in $\mu$ whose images under $T$ fall into $B$.}
\end{highlight}

Intuitively:
\begin{itemize}
    \item If $x$ is sampled from $\mu$, then $T(x)$ is sampled from $T_\#\mu$.
    \item So, $T_\#\mu$ is the \textbf{distribution of the image} of $x$ under $T$.
\end{itemize}

\section*{3. A Simple Illustration}

Let us build some intuition with a concrete example:

\begin{itemize}
    \item Take $X = \mathbb{R}$ and let $\mu$ be the uniform distribution on the interval $[0,1]$.
    \item Define a map $T : [0,1] \to [0,1]$ by $T(x) = 1 - x$.
\end{itemize}

What happens?
\begin{itemize}
    \item Every point $x$ is reflected around $x = 0.5$.
    \item For example: $T(0.2) = 0.8$, $T(0.5) = 0.5$, $T(0.9) = 0.1$.
    \item The pushforward $T_\#\mu$ is still uniform on $[0,1]$, but the ``orientation'' has flipped.
\end{itemize}

\begin{highlight}
This illustrates that a pushforward can change the structure or layout of a distribution without changing its total mass or support.
\end{highlight}

This kind of reflection is a trivial example, but the same idea generalizes: by applying the right transformation $T$, we can morph one distribution into another in a principled way.

\section*{4. The Big Goal of Optimal Transport}

Suppose we are given two distributions:
\begin{itemize}
    \item $\mu$ on a space $X$
    \item $\nu$ on a space $Y$
\end{itemize}

The goal of optimal transport is:

\begin{highlight}
\textbf{Find a map $T : X \to Y$ such that $T_\#\mu = \nu$, and the overall cost of moving each $x$ to $T(x)$ is as small as possible.}
\end{highlight}

The ``cost'' is defined via a function $c(x, y)$, often $\|x - y\|$ or $\|x - y\|^2$.
Different choices of cost lead to different flavors of OT (e.g., Wasserstein-1, Wasserstein-2).

\vspace{1em}
\begin{center}
\textit{Understanding probability measures and pushforwards is the essential starting point for mastering the rest of Optimal Transport theory.}
\end{center}


\newpage
\lhead{Level 1}
\section*{Level 1 – Monge, Kantorovich, and Duality}

\subsection*{1. The Monge Formulation (1781)}

Gaspard Monge proposed a very intuitive version of the transport problem:

\begin{highlight}
Given two distributions $\mu$ (source) and $\nu$ (target), find a transport map $T : X \to Y$ such that $T_\#\mu = \nu$ and minimizes the total cost
\[
\int_X c(x, T(x))\,d\mu(x).
\]
\end{highlight}

\noindent Here, $c(x, y)$ is the \textbf{cost of moving} one unit of mass from $x$ to $y$.

\textbf{Key limitation:} Monge’s formulation does not allow \emph{splitting mass}. If a point mass must go somewhere, it can only go to one destination.

\vspace{0.5em}
This becomes problematic, for instance, if:
\begin{itemize}
  \item $\mu$ is a single point (Dirac delta), and
  \item $\nu$ is spread out.
\end{itemize}
Then no such $T$ exists.

\subsection*{2. The Kantorovich Relaxation (1942)}

Leonid Kantorovich proposed a relaxation that allows mass to split.

\begin{highlight}
Instead of finding a transport map $T$, find a joint probability measure $\pi(x, y)$ over $X \times Y$ with marginals $\mu$ and $\nu$.
\end{highlight}

This $\pi$ tells us how much mass is transported from $x$ to $y$.

\textbf{Formal problem:}
\[
\inf_{\pi \in \Pi(\mu, \nu)} \int_{X \times Y} c(x, y)\,d\pi(x, y)
\]
where $\Pi(\mu, \nu)$ is the set of all couplings with marginals $\mu$ and $\nu$.

\vspace{0.5em}
This formulation:
\begin{itemize}
  \item Always has a solution under mild conditions
  \item Can be efficiently optimized
  \item Includes Monge’s case as a special (deterministic) scenario
\end{itemize}

\section*{From Monge to Kantorovich: Why the Relaxation Is Needed}

\subsection*{The Problem with Monge}
In Monge's formulation, we try to find a \emph{map} $T : X \to Y$ that pushes one distribution $\mu$ onto another $\nu$, minimizing a transport cost:
\[
\inf_{T : T_\#\mu = \nu} \int_X c(x, T(x)) \, d\mu(x)
\]

But this approach has a major limitation:

\begin{highlight}
It does not allow \textbf{splitting mass}. Each point $x$ must be sent entirely to a single $T(x)$.
\end{highlight}

This becomes an issue when, for instance:
\begin{itemize}
  \item $\mu$ is a single Dirac mass (e.g., $\delta_x$)
  \item $\nu$ is a distribution over multiple points (e.g., a mixture of Diracs)
\end{itemize}

In such cases, there's \emph{no way} to find a map $T$ such that $T_\#\mu = \nu$. The Monge problem has \textbf{no solution}.

\subsection*{The Kantorovich Idea: A Probabilistic View}
To solve this, Kantorovich proposed a \textbf{relaxation}:

\begin{highlight}
Rather than assigning a single destination $T(x)$ for each $x$, assign a \emph{distribution over possible destinations}. That is: a \textbf{coupling} $\pi(x, y)$.
\end{highlight}

Formally:
\begin{itemize}
  \item Let $\pi \in \Pi(\mu, \nu)$ be a joint probability distribution over $X \times Y$.
  \item The marginals must satisfy:
  \[
  \int_Y \pi(x, y)\,dy = \mu(x), \quad \int_X \pi(x, y)\,dx = \nu(y)
  \]
  \item The cost becomes:
  \[
  \inf_{\pi \in \Pi(\mu, \nu)} \int_{X \times Y} c(x, y)\, d\pi(x, y)
  \]
\end{itemize}

Now:
\begin{itemize}
  \item Mass at $x$ can be split across multiple $y$ destinations.
  \item The problem is \textbf{convex} and always has a solution under mild conditions.
  \item When a deterministic map exists, the solution $\pi$ collapses to a Monge map: $\pi = (\text{id}, T)_\#\mu$.
\end{itemize}

\subsection*{Summary}
\begin{center}
\begin{tcolorbox}[colback=white,colframe=blue!75!black,title=Monge vs. Kantorovich]
\begin{itemize}
  \item \textbf{Monge:} deterministic, no splitting, can fail to exist
  \item \textbf{Kantorovich:} probabilistic, allows splitting, always solvable
  \item \textbf{Key bridge:} when a Monge map exists, it's embedded in Kantorovich’s solution
\end{itemize}
\end{tcolorbox}
\end{center}

This relaxation is what opened the door to a full mathematical theory of transport---from convex analysis to duality, geometry, and machine learning applications.


\subsection*{3. Duality: The Power of Potentials}

The Kantorovich problem has a beautiful dual form:

\[
\sup_{u, v \in \mathcal{C}} \left[ \int u(x)\,d\mu(x) + \int v(y)\,d\nu(y) \right] \quad \text{such that } u(x) + v(y) \leq c(x, y)
\]

Here:
\begin{itemize}
  \item $u$ and $v$ are called \textbf{Kantorovich potentials}
  \item They encode the optimal “price” structure of transport
\end{itemize}

Duality connects optimization to geometry. It’s the foundation of modern computational OT and allows recovery of the optimal transport map in special cases (e.g., for quadratic costs).

\section*{From Kantorovich to the Dual Form}

\subsection*{Why Do We Introduce a Dual?}
In optimization, the \textbf{dual problem} offers a different perspective on the same question. Often:
\begin{itemize}
  \item The primal formulation describes \emph{what} we want to do.
  \item The dual tells us \emph{how costly} the constraints are.
  \item Solving the dual can be easier, and yields important insights into the structure of the solution.
\end{itemize}

To build a dual, we use a powerful tool from constrained optimization: the \textbf{Lagrangian}. It combines the original objective with its constraints using \emph{multipliers} that measure how much each constraint “hurts.”

\subsection*{Step-by-Step: The Lagrangian Approach}
We start with the \textbf{primal Kantorovich problem}:
\[
\inf_{\pi \in \Pi(\mu, \nu)} \int_{X \times Y} c(x, y) \, d\pi(x, y)
\]
subject to:
\[
\int_Y \pi(x, y)\, dy = \mu(x), \quad \int_X \pi(x, y)\, dx = \nu(y)
\]

To incorporate the constraints, we use the method of \textbf{Lagrange multipliers}. Let:
\begin{itemize}
  \item $u : X \to \mathbb{R}$ be a Lagrange multiplier (or "potential") for the marginal constraint on $\mu$
  \item $v : Y \to \mathbb{R}$ for the marginal on $\nu$
\end{itemize}

Introduce Lagrange multipliers:
\begin{itemize}
  \item $u(x)$ for the constraint $\int_Y \pi(x, y)\,dy = \mu(x)$
  \item $v(y)$ for the constraint $\int_X \pi(x, y)\,dx = \nu(y)$
\end{itemize}
We can think of $u(x)$ and $v(y)$ as the \textbf{costs} associated with moving mass from $x$ to $y$.
The Lagrangian $\mathcal{L}$ combines the primal objective with the constraints:

\[
\mathcal{L}(\pi, u, v) = \int_{X \times Y} c(x, y) \, d\pi(x, y)
- \int_X u(x)\left[\int_Y d\pi(x, y) - d\mu(x)\right]
- \int_Y v(y)\left[\int_X d\pi(x, y) - d\nu(y)\right]
\]
Reorganizing:
\begin{align*}
\mathcal{L}(\pi, u, v) =
&\int_{X \times Y} \left[ c(x, y) - u(x) - v(y) \right] d\pi(x, y) \\
&+ \int_X u(x)\,d\mu(x) + \int_Y v(y)\,d\nu(y)
\end{align*}

Now, take the \textbf{infimum over all couplings} $\pi$:
\begin{itemize}
  \item If $c(x, y) - u(x) - v(y) < 0$ anywhere, the infimum is $-\infty$.
  \item So to keep $\mathcal{L}$ bounded below, we must enforce:
\[
\boxed{u(x) + v(y) \leq c(x, y)} \quad \text{for all } (x, y)
\]
\end{itemize}

Hence, the \textbf{dual problem} becomes:
\[
\sup_{u, v} \left\{ \int_X u(x)\, d\mu(x) + \int_Y v(y)\, d\nu(y) \right\} \quad \text{subject to } u(x) + v(y) \leq c(x, y)
\]

This is the \textbf{Kantorovich dual formulation}.

\vspace{1em}
\begin{tcolorbox}[colback=blue!5!white, colframe=blue!75!black, title=Takeaway]
The dual variables $u$ and $v$ can be interpreted as \emph{economic potentials}---they represent how much we're willing to pay to move mass from $x$ to $y$, given cost $c(x, y)$.
\end{tcolorbox}




\subsection*{4. Summary Table}

\begin{center}
\begin{tabular}{|c|c|c|}
\hline
\textbf{Formulation} & \textbf{Object} & \textbf{Allows Splitting?} \\\hline
Monge & $T : X \to Y$ & No \\\hline
Kantorovich (Primal) & $\pi(x, y)$ & Yes \\\hline
Kantorovich (Dual) & $u(x), v(y)$ & Yes \\\hline
\end{tabular}
\end{center}

\vspace{1em}
\begin{center}
\textit{Monge is beautiful but rigid. Kantorovich is flexible and powerful. Duality is the lens that unites them.}
\end{center}

\newpage
\lhead{Level 2}
\section*{Level 2 - Understanding Wasserstein Distances}

\subsection*{The Big Idea: Distance Between Distributions}
Imagine you have two piles of sand: one represents $\mu$, the other $\nu$. Wasserstein distance answers:

\begin{highlight}
\textbf{How much effort does it take to reshape $\mu$ into $\nu$?}
\end{highlight}

But unlike most distances between distributions (e.g., KL divergence, Total Variation), Wasserstein is \emph{geometry-aware}. It respects the actual locations of the mass.

\subsection*{The Formal Definition (Wasserstein-$p$)}
Let $p \geq 1$ and $c(x,y) = \|x - y\|^p$. Then the \textbf{Wasserstein-$p$ distance} is:
\[
W_p(\mu, \nu) = \left( \inf_{\pi \in \Pi(\mu, \nu)} \int_{X \times Y} \|x - y\|^p \, d\pi(x, y) \right)^{1/p}
\]

Key components:
\begin{itemize}
  \item $\Pi(\mu, \nu)$: all joint distributions (couplings) with marginals $\mu$ and $\nu$
  \item $\|x - y\|^p$: cost of moving a unit mass from $x$ to $y$
  \item The $p$-th root makes it a proper distance
\end{itemize}

\subsection*{Interpreting $W_p$}
\begin{itemize}
  \item $W_1$: Total work — move mass over minimum distance (used in WGAN)
  \item $W_2$: Minimizes squared distances — behaves smoothly, useful in theory
  \item $W_\infty$: The farthest distance any mass must travel
\end{itemize}

\begin{tcolorbox}[colback=white, colframe=blue!50!black, title=Geometric Meaning]
\textbf{Wasserstein distance} measures the ``optimal way'' to morph one shape into another — mass by mass — according to physical cost.
\end{tcolorbox}

\subsection*{Why It Matters}
\begin{itemize}
  \item Wasserstein distances define \emph{meaningful} geometry over the space of probability distributions
  \item They make sense even when distributions have disjoint support (unlike KL)
  \item They are the foundation of generative models, image restoration, domain adaptation, and more
\end{itemize}

\vspace{0.5em}
\textbf{Example:} Let $\mu$ be the uniform distribution on the interval $[0, 1]$, and let $\nu$ be uniform on $[2, 3]$. That is:
\[
\mu = \mathcal{U}[0, 1], \quad \nu = \mathcal{U}[2, 3]
\]

The mass is the same in both distributions — but it's located in different places.

Then:
\[
W_1(\mu, \nu) = \int_0^1 |x - (x + 2)|\,dx = \int_0^1 2\,dx = 2
\]
\[
W_2(\mu, \nu)^2 = \int_0^1 (x - (x + 2))^2 \, dx = \int_0^1 4 \, dx = 4 \quad \Rightarrow \quad W_2(\mu, \nu) = 2
\]

\begin{center}
\textit{The optimal transport map here is $T(x) = x + 2$: a pure shift.}\\
Each mass element $x$ from $[0,1]$ is simply moved to $x+2 \in [2,3]$.
\end{center}

\subsubsection*{Example: Asymmetric Mass Shift}

Let:
\[
\mu = 0.9\,\delta_0 + 0.1\,\delta_{10}, \quad \nu = \delta_1
\]

So:
\begin{itemize}
  \item 90\% of the mass in $\mu$ is at $x = 0$,
  \item 10\% is far away, at $x = 10$,
  \item $\nu$ has all its mass at $y = 1$.
\end{itemize}

We compute the optimal transport:

\begin{itemize}
  \item Move $0.9$ mass from $0 \to 1$
  \item Move $0.1$ mass from $10 \to 1$
\end{itemize}

Then:

\[
W_1(\mu, \nu) = 0.9 \cdot |0 - 1| + 0.1 \cdot |10 - 1| = 0.9 \cdot 1 + 0.1 \cdot 9 = 0.9 + 0.9 = 1.8
\]

\[
W_2^2(\mu, \nu) = 0.9 \cdot (0 - 1)^2 + 0.1 \cdot (10 - 1)^2 = 0.9 \cdot 1 + 0.1 \cdot 81 = 0.9 + 8.1 = 9.0
\quad \Rightarrow \quad W_2(\mu, \nu) = 3
\]

\begin{tcolorbox}[colback=blue!2!white, colframe=blue!90!black, title=Conclusion]
Here, $W_1 = 1.8$, but $W_2 = 3$ — the outlier at $x = 10$ increases $W_2$ much more than $W_1$.
\end{tcolorbox}

\begin{center}
\textit{Wasserstein-$1$ sees “how far things move”.\\
Wasserstein-$2$ sees “how badly they move”.}
\end{center}

\newpage
\lhead{Level 3}
\section*{Level 3 – Optimal Transport in Generative Models}

\subsection*{Why Use Optimal Transport in Generative Models?}

When we build a generative model, we want to turn simple random noise (like a Gaussian) into complex data — like images, speech, or molecules. 

\begin{highlight}
\textbf{Optimal Transport (OT) tells us how to move one distribution into another — in the most efficient way.}
\end{highlight}

In this setting:
\begin{itemize}
  \item $\mu$ is our known distribution (e.g., standard normal noise),
  \item $\nu$ is the target data distribution (e.g., real images),
  \item The OT map $T^*$ gives us the best way to turn samples from $\mu$ into samples from $\nu$.
\end{itemize}

\subsection*{Two Strategies}

\begin{enumerate}
  \item \textbf{Use OT distance as a loss function:} like in Wasserstein GANs (WGAN), we use $W_1(\mu, \nu)$ to guide the generator $G$.
  \item \textbf{Use the OT map directly as the generator:} the idea in the paper. Here, we don’t need a discriminator — we directly learn the map $T^*$.
\end{enumerate}

\subsection*{Why $W_2$ Instead of $W_1$?}

The paper focuses on $W_2$ — the squared Wasserstein distance — because it has a powerful structure:

\begin{itemize}
  \item Under mild conditions, the optimal map has the form:
  \[
  T^*(x) = \nabla \psi(x), \quad \text{for a convex function } \psi.
  \]
  \item This means we can learn $\psi$ as a neural network, and recover $T$ just by taking its gradient.
\end{itemize}

This would not be possible with $W_1$ — we’d get a distance, but not an explicit map.

\begin{tcolorbox}[colback=blue!4!white, colframe=blue!80!black, title=Bottom Line]
OT gives us more than a number — it gives us a generator.
With $W_2$, we can turn random noise into structured data without adversarial training.
\end{tcolorbox}

\subsection*{What Happens Next}
In the next levels, we’ll see:
\begin{itemize}
  \item The strength of $W_2$ for deep theoretical and practical reasons,
  \item How to learn $\psi$ and $T$ from data,
  \item And how this builds real-world generative models (e.g., for image synthesis, restoration).
\end{itemize}

\newpage
\lhead{Level 4}
\section*{Level 4 – Learning the Optimal Transport Map}

\subsection*{Objective: Not Just Distance — the Map Itself}
In many applications, we are not just interested in the value of the Wasserstein distance, but in the \textbf{optimal map} $T^*$ that pushes $\mu$ to $\nu$.

\begin{highlight}
\textbf{The goal is to learn a function $T$ such that $T_\#\mu = \nu$ and the transport cost $\int c(x, T(x)) \, d\mu(x)$ is minimized.}
\end{highlight}

This leads us into the realm of \textbf{map-based optimal transport}. Here, the map $T$ can be seen as a 

\begin{itemize}
  \item \textbf{Generator:} in generative models,
  \item \textbf{Restorer:} in image restoration tasks,
  \item \textbf{Interpolator:} in shape morphing or style transfer.
\end{itemize}

\subsection*{From Dual to Map: Brenier's Insight (W$_2$ Case)}
When the cost is quadratic $c(x,y) = \frac{1}{2}\|x - y\|^2$, and $\mu$ is absolutely continuous:

\begin{tcolorbox}[colback=blue!5!white, colframe=blue!75!black, title=Brenier's Theorem (1991)]
There exists a convex function $\phi : \mathbb{R}^d \to \mathbb{R}$ such that the optimal map is:
\[
T^*(x) = \nabla \phi(x)
\]
\end{tcolorbox}

This result is foundational — it means that \textbf{the optimal transport map is a gradient of a convex potential}. This links OT to convex analysis, PDEs, and deep learning.

\subsection*{Learning $T$ via a Saddle Point Formulation}
In practice, we often solve a saddle point problem over two neural networks:
\begin{itemize}
  \item $\psi(y)$ approximates a convex potential
  \item $T(x)$ is trained to match the gradient of $\psi$
\end{itemize}

An example (for $W_2$):
\[
\min_{\psi} \max_{T} \left[ \mathbb{E}_{x \sim \mu} \langle x, T(x) \rangle - \psi(T(x)) + \mathbb{E}_{y \sim \nu} \psi(y) \right]
\]

This yields $T^*(x) \approx \nabla \psi^*(x)$, the gradient of the convex conjugate, when $\psi$ is optimal.

\subsection*{In Detail: The Special Case of Quadratic Cost, $W_2$}
Let the cost be $c(x,y) = \frac{1}{2}\|x - y\|^2$. The primal Kantorovich formulation becomes:
\[
\frac{1}{2} W_2^2(\mu, \nu) = \inf_{\pi \in \Pi(\mu, \nu)} \int \frac{1}{2}\|x - y\|^2 \, d\pi(x,y)
\]

The dual formulation is:
\[
\sup_{u, v} \left\{ \int u(x)\,d\mu(x) + \int v(y)\,d\nu(y) \right\} \quad \text{subject to } u(x) + v(y) \leq \frac{1}{2}\|x - y\|^2
\]

Now introduce the \textbf{$c$-transform}. For a function $v$, define its $c$-transform:
\[
v^c(x) := \inf_{y \in \mathbb{R}^d} \left\{ \frac{1}{2}\|x - y\|^2 - v(y) \right\}
\]

Similarly, $u^c(y) := \inf_{x} \left\{ \frac{1}{2}\|x - y\|^2 - u(x) \right\}$.

\subsection*{Reduction to a Single Potential via Legendre Transform}
Set:
\[
\psi(y) := \frac{1}{2}\|y\|^2 - v(y)
\quad \Rightarrow \quad v(y) = \frac{1}{2}\|y\|^2 - \psi(y)
\]
Then the $c$-transform becomes:
\begin{align*}
v^c(x)
&= \inf_y \left\{ \frac{1}{2}\|x - y\|^2 - \left( \frac{1}{2}\|y\|^2 - \psi(y) \right) \right\} \\
&= \inf_y \left\{ \frac{1}{2}\|x\|^2 - \langle x, y \rangle + \psi(y) \right\} \\
&= \frac{1}{2}\|x\|^2 - \sup_y \left\{ \langle x, y \rangle - \psi(y) \right\} \\
&= \frac{1}{2}\|x\|^2 - \psi^*(x)
\end{align*}
where $\psi^*$ is the \textbf{Legendre (convex) conjugate} of $\psi$.

The dual becomes:
\[
\frac{1}{2} W_2^2(\mu, \nu) = \int \frac{1}{2}\|x\|^2 \, d\mu(x) + \int \frac{1}{2}\|y\|^2 \, d\nu(y) - \inf_{\psi \text{ convex}} \left\{ \int \psi^*(x) \, d\mu(x) + \int \psi(y) \, d\nu(y) \right\}
\]

This is a fundamental variational problem in convex analysis. With $\psi$ optimal, the optimal map is the gradient of its convex conjugate (Brenier's potential is $\phi = \psi^*$):
\[
T^*(x) = \nabla \psi^*(x)
\]

\subsection*{Learning $\psi$ via Neural Networks}
In practice, we parameterize $\psi$ by a neural network and optimize:
\[
\min_\psi \left\{ \mathbb{E}_{x \sim \mu}[\psi^*(x)] + \mathbb{E}_{y \sim \nu}[\psi(y)] \right\}
\]

The difficulty is $\psi^*$: evaluating the conjugate is itself an optimisation problem for every $x$. This is what motivates the saddle-point reformulation of Level 5, which yields both the OT distance and the optimal transport map.

\subsection*{Applications and Benefits}
\begin{itemize}
  \item Learn a \textbf{generator} $T = \nabla \psi^*$ that is structure-aware and cost-efficient.
  \item \textbf{Better stability} than adversarial GANs: the inner maximisation has a clear meaning (a conjugate), and there is no discriminator to collapse.
  \item \textbf{Interpretability}: the potential $\psi$ encodes a landscape that shapes how mass moves.
\end{itemize}

\begin{tcolorbox}[colback=blue!3!white, colframe=blue!80!black, title=Summary]
When $c(x,y) = \frac{1}{2}\|x - y\|^2$, OT reduces to a convex variational problem with a single potential function.
\textbf{The optimal map is the gradient of the convex conjugate of that potential.}
\end{tcolorbox}



\newpage
\lhead{Level 5}
\section*{Level 5 – Practical Recovery of $\psi$ and $T$ in the Paper}

In the ICLR 2022 paper \textit{Generative Modeling with Optimal Transport Maps}, the authors propose a practical method to recover both the potential $\psi$ and the optimal transport map $T$ for the quadratic cost (Wasserstein-2).

\subsection*{From Theory to Practice: Saddle Point Formulation}

We begin from the dual reformulated with a single convex potential $\psi$:
\[
\inf_{\psi} \left\{ \mathbb{E}_{x \sim \mu} \left[ \psi^*(x) \right] + \mathbb{E}_{y \sim \nu} \left[ \psi(y) \right] \right\}
\]
Instead of computing the Legendre conjugate $\psi^*$ explicitly, the authors propose a saddle-point reformulation:
\[
\inf_{\psi} \sup_{T} \left\{ \mathbb{E}_{x \sim \mu} \left[ \langle x, T(x) \rangle - \psi(T(x)) \right] + \mathbb{E}_{y \sim \nu} \left[ \psi(y) \right] \right\}
\]

This allows simultaneous learning of $\psi$ and $T$.

\subsection*{Optimization Procedure (Algorithm 1)}

\begin{itemize}
  \item $\psi$ is parameterized by a neural network $\psi_\omega : \mathbb{R}^D \to \mathbb{R}$.
  \item $T$ is also parameterized by a neural network $T_\theta : \mathbb{R}^D \to \mathbb{R}^D$.
  \item The authors perform alternating stochastic gradient steps:
  \begin{itemize}
    \item \textbf{Step 1:} Fix $\psi$, maximize over $T$:
    \[
    \mathbb{E}_{x \sim \mu} \left[ \langle x, T(x) \rangle - \psi(T(x)) \right]
    \]
    \item \textbf{Step 2:} Fix $T$, minimize over $\psi$:
    \[
    \mathbb{E}_{x \sim \mu} \left[ -\psi(T(x)) \right] + \mathbb{E}_{y \sim \nu} \left[ \psi(y) \right]
    \]
  \end{itemize}
  \item This is done using mini-batches drawn from $\mu$ and $\nu$, with Adam optimizer.
  \item The update frequencies $K_G$ (for $T$) and $K_\psi$ (for $\psi$) can differ (e.g. $K_G = 5$, $K_\psi = 1$).
\end{itemize}

\subsection*{Theoretical Justification}
If $\psi$ is convex and smooth, then:
\[
T^*(x) = \nabla \psi^*(x) \quad \text{and} \quad \nabla \psi(T^*(x)) = x
\]

Thus, $T$ approximates the \emph{inverse} of $\nabla \psi$. The equilibrium of the saddle point problem yields both:
\begin{itemize}
  \item The optimal transport cost $W_2^2(\mu, \nu)$
  \item The optimal map $T^*(x)$ without explicitly computing $\psi^*$
\end{itemize}

\subsection*{Stabilization: Regularization and Penalties}
To stabilize training, the authors add a regularization term to the $\psi$-update, such as:
\[
\lambda \cdot \left\| \mathbb{E}_{x \sim \mu} \nabla \psi(T(x)) - x \right\|^2
\]
which vanishes at the optimum due to the theoretical condition $\nabla \psi(T^*(x)) = x$.

This regularizer improves convergence and accuracy, especially in high-dimensional settings.

\subsection*{Summary of What Is Learned}
\begin{itemize}
  \item $\psi$: a convex function (neural network) modeling the potential.
  \item $T$: a transport map, trained to maximize the dual objective.
\end{itemize}

The final result is a map $T$ that can be directly applied as a generator in high-dimensional image space — with no need for autoencoders or adversarial critics.

\begin{tcolorbox}[colback=blue!2!white, colframe=blue!80!black, title=In Practice]
The paper minimizes $\psi$ and maximizes $T$ alternately — solving a saddle point problem that recovers both the OT map and the potential.
\end{tcolorbox}


\newpage
\lhead{Level 6}
\section*{Level 6 – Generalization to Unequal Dimensions}

\subsection*{The Challenge}

In many practical generative tasks, we want to map from a low-dimensional latent space $\mathbb{R}^H$ (e.g., $H = 100$) to a high-dimensional data space $\mathbb{R}^D$ (e.g., $D = 784$ for images). 

But classical OT theory assumes $\mu$ and $\nu$ lie in the \emph{same} space.

\begin{highlight}
\textbf{How can we define optimal transport when $H \neq D$?}
\end{highlight}

\subsection*{The Solution: Embed the Latent Space}

The paper resolves this by introducing a fixed embedding function:
\[
Q : \mathbb{R}^H \to \mathbb{R}^D
\]
This map projects the low-dimensional space into the same ambient space as the data.

Once embedded, we define the cost:
\[
c(x, y) = \frac{1}{2} \|Q(x) - y\|^2
\]
and solve the OT problem in $\mathbb{R}^D$.

\subsection*{The Min-Max Formulation Still Works}

With this new cost, the dual objective becomes:
\[
\min_{\psi} \max_{G} \ 
\mathbb{E}_{x \sim \mu} \left[ \langle Q(x), G(x) \rangle - \psi(G(x)) \right] + 
\mathbb{E}_{y \sim \nu} \left[ \psi(y) \right]
\]

Here:
\begin{itemize}
  \item $G : \mathbb{R}^H \to \mathbb{R}^D$ is the generator,
  \item $\psi : \mathbb{R}^D \to \mathbb{R}$ is the potential function,
  \item $Q$ is a fixed upsampling map (e.g., bicubic interpolation from $8\times8$ to $32\times32$).
\end{itemize}

\subsection*{Why This Is Important}

\begin{itemize}
  \item It allows the method to work with latent variables — just like GANs or VAEs.
  \item It bypasses the need for an encoder: we only go from noise to data, not back.
  \item It enables direct generation in image space, using only optimal transport.
\end{itemize}

\begin{tcolorbox}[colback=blue!5!white, colframe=blue!80!black, title=Key Idea]
By embedding $\mu$ into the ambient space using $Q$, we can define a valid OT problem even when dimensions differ.
The learned generator $G$ replaces the transport map.
\end{tcolorbox}

\subsection*{In Practice}
In the experiments:
\begin{itemize}
  \item $x \sim \mu$ is sampled from $\mathcal{N}(0, I_H)$,
  \item $Q(x)$ is a low-res upsampling (e.g., from $8\times8$ to $32\times32$),
  \item $G(x)$ learns to refine this into a realistic sample from $\nu$.
\end{itemize}

This approach preserves the geometry of OT while maintaining compatibility with standard generative modeling pipelines.

\vspace{1em}

\lhead{Level 7}
\section*{Level 7 – Experiments and Limitations}

\subsection*{Generative Modeling from Noise}

The authors evaluate their method on classical generative tasks:
\begin{itemize}
  \item \textbf{Datasets:} MNIST ($28\times28$), CIFAR10 ($32\times32$), CelebA ($64\times64$).
  \item \textbf{Latent distribution:} $x \sim \mathcal{N}(0, I)$ in $\mathbb{R}^H$.
  \item \textbf{Fixed embedding $Q$:} Bicubic upsampling of $8\times8$ latent noise to full resolution.
\end{itemize}

The generator $G$ learns to push these embedded samples toward real data.

\subsection*{Performance (FID Score)}

Fréchet Inception Distance (FID) is used to measure quality of generated samples. Lower is better.

\begin{itemize}
  \item On MNIST and CIFAR10, the OT-based method matches or improves upon WGAN, AE-OT, and OT-GAN.
  \item On CelebA, results are visually sharp and FID is competitive.
\end{itemize}

\subsection*{Unpaired Image Restoration}

The model is also tested on tasks where noisy inputs $\mu$ must be mapped to clean targets $\nu$:
\begin{itemize}
  \item \textbf{Denoising:} Gaussian noise added to images.
  \item \textbf{Inpainting:} Random patches are masked.
  \item \textbf{Colorization:} Grayscale to color images.
\end{itemize}

Despite being unpaired (no $x \leftrightarrow y$ examples), the learned map $G$ recovers realistic outputs.

\subsection*{Limitations and Open Questions}

\begin{itemize}
  \item \textbf{Scalability:} Although effective at $32\times32$ or $64\times64$, the model has not yet been scaled to high-res (e.g., $256\times256$).
  \item \textbf{Choice of $Q$:} The fixed embedding $Q$ works well empirically, but its design is somewhat heuristic (e.g., bicubic). Learning $Q$ jointly might improve flexibility.
  \item \textbf{Convexity Assumption:} The method relies on approximating $\psi$ as convex. While Input Convex Neural Networks (ICNNs) can enforce this, they limit architectural choices.
  \item \textbf{Optimization Stability:} The saddle-point formulation is simpler than GANs but still prone to local optima and training instabilities.
\end{itemize}

\begin{tcolorbox}[colback=blue!5!white, colframe=blue!80!black, title=Key Insight]
OT maps can be learned directly in image space — without adversarial losses, and with minimal supervision. 
The method scales well to small- and mid-resolution domains, while leaving room for extensions.
\end{tcolorbox}


\end{document}
